% ==============================================================================
% PDF-1 / Section 02 — Phase Space Construction (Skeleton)
% No new primitives. Diagnostic-only. Container structure.
% ==============================================================================

\section{Phase Space Construction}
\label{sec:phase-space}

\subsection{Objective of the construction}
This section defines the \emph{structural container} for constructing a diagnostic
phase space over the Enterprise Continuum \(K_{EA}\).
The goal is to provide a clear slot for mapping observable enterprise states into
a formally bounded space without introducing new variables, metrics, or primitives.

All coordinates, axes, and constraints referenced here are assumed to be inherited
from KEA (PDF-0) and \texttt{ETS.K\_EA.v1.1}. No extensions are permitted in this run.

\subsection{State representation (container)}
An enterprise state is represented as a point (or region) in the admissible
phase space defined by:
\begin{itemize}
  \item the existing state descriptors of \(K_{EA}\),
  \item the admissible configuration set as specified by the ETS,
  \item explicit boundary conditions inherited from the STOP logic.
\end{itemize}

This subsection intentionally omits any numerical parametrization or computation.
It exists solely to anchor later executable demonstrations to a stable formal slot.

\subsection{Admissible regions and boundaries}
The phase space is partitioned into admissible and non-admissible regions according
to constraints already defined in the immutable sources.
Boundary surfaces are treated as diagnostic objects (not optimization targets).

No assumptions are made about smoothness, convexity, or continuity beyond what is
explicitly stated in the underlying KEA specification.

\subsection{Observability and projection}
Observed enterprise signals are projected into the phase space via diagnostic
projection operators already defined elsewhere.
This document does not define new projection rules; it only reserves the structural
location where such projections are referenced and discussed.

\subsection{Non-goals}
\begin{itemize}
  \item No construction of new dimensions or axes.
  \item No optimization criteria or target states.
  \item No intervention logic or control objectives.
  \item No claims about completeness or sufficiency of observation.
\end{itemize}
